\documentclass[14pt,a4paper]{article}
\usepackage[a4paper, % размер страницы,
text={180mm, 260mm},% ширина и высота текста,
left=15mm, top=15mm]{geometry} % размеры пустых полей,
\usepackage[utf8]{inputenc} %кодировка текста,
\usepackage{cmap}  %кодировка pdf-документа,
\usepackage[english,russian]{babel} % поддержка русского языка,
\usepackage{indentfirst} % корректировка отступов,
\usepackage{amssymb} % математические символы,
\usepackage{amsmath} % математические конструкции,
\usepackage{graphicx}

\usepackage[backend=biber,style=gost-numeric]{biblatex}
\usepackage{xcolor}

\addbibresource{D:/Rozhkova/Projects/DL/LaTex/bibliography.bib}

\begin{document}

\begin{abstract}
В промышленном комплексе Республики Беларусь контроль качества продукции является одним из ключевых факторов конкурентоспособности на внутреннем и внешнем рынках. В условиях реализации стратегий устойчивого развития, курса на импортозамещение и цифровую трансформацию экономики особую значимость приобретает внедрение автоматизированных систем контроля качества. Автоматизированные средства контроля способны функционировать в режиме реального времени, оперативно реагировать на возникновение отклонений и тем самым повышать эффективность производственных линий.

Среди отраслей, где визуальный контроль качества является критически важным, можно выделить металлургическое производство, железнодорожный транспорт, диагностику линий электропередач, а также широкий спектр промышленных технологических процессов. В последние годы особое развитие получили исследования, посвящённые применению автоматизированного визуального контроля в гражданском строительстве и мониторинге технического состояния инженерных сооружений. К таким задачам относятся обнаружение трещин в бетонных конструкциях, автоматический контроль состояния дорожных покрытий и тротуаров, выявление выбоин, шероховатостей и других поверхностных дефектов, которые могут негативно сказаться на безопасности эксплуатации объектов . Подобные подходы демонстрируют высокую актуальность в связи с необходимостью своевременного проведения ремонтно-восстановительных работ и оптимизации затрат на обслуживание инфраструктуры.

Методы обнаружения дефектов находят применение не только в производстве промышленных изделий — таких как металлические заготовки, текстильные полотна, полупроводниковые пластины или композиционные материалы, — но и в агропромышленном секторе, где требуется своевременная сортировка сельскохозяйственной продукции и выявление её повреждений. Кроме того, эти подходы востребованы при обследовании строительных объектов, мостов, дорожных полотен и других элементов городской инфраструктуры, включая бетонные и асфальтовые покрытия. Таким образом, задачи автоматизированного обнаружения аномалий и дефектов носят межотраслевой характер и являются важным направлением развития современных интеллектуальных систем визуального анализа.

В последнее время глубокое обучение доказало свою исключительную эффективность в таких задачах, как обнаружение и классификация объектов, распознавание лиц, идентификация образов, диагностика неисправностей, отслеживание целей и множестве других приложений, связанных с обработкой изображений. Модели глубоких нейронных сетей демонстрируют высокую устойчивость к изменению условий съёмки — фону, освещённости, цветовой гаммы, формы и размеров объектов, варьирующей интенсивности изображения. Эти свойства особенно важны при обнаружении сложных поверхностных дефектов, характерных для промышленных условий, где изображения могут быть зашумлены, иметь низкое качество или быть сняты при непостоянных параметрах окружающей среды.  Помимо простого выявления наличия дефекта, современные системы нередко требуют точной оценки его параметров: размеров, формы, протяжённости и типа. Такие возможности особенно востребованы в задачах визуального контроля качества продукции и автоматизации производственных процессов.

Целями настоящего обзора являются анализ существующих подходов к решению задач обнаружения аномалий на изображениях с использованием методов глубокого обучения, выявление их преимуществ и ограничений, а также обоснование выбора методик и формулировка целей собственного исследования в рамках подготовки магистерской диссертации. Проведённый обзор позволит определить наиболее перспективные направления развития алгоритмов визуального контроля и сформировать основу для построения собственной экспериментальной методологии.
\end{abstract}

\section{Введение}
\subsection*{Определения поверхностных дефектов}

В работе \cite{Hutten2024} поверхностные дефекты определяются как участки с локальными отклонениями текстуры, цвета или структуры, не соответствующие нормальному фону поверхности. В промышленности выделяют два основных типа таких дефектов, связанных с физическими свойствами материалов:
\begin{enumerate}
	\item Дефекты на поверхностях жёстких материалов (металл, дерево, керамика). Визуально проявляются в виде нерегулярных трещин, сколов, включений, царапин и аналогичных повреждений.

	\item Дефекты на поверхностях гибких материалов (ткань, кожа, резина). К ним относятся пузыри, точечные дефекты, а также короткие линейные дефекты.
\end{enumerate}
Независимо от типа, поверхностные дефекты влияют не только на внешний вид изделия, но и могут стать начальной стадией развития механических повреждений или химической коррозии, что снижает качество и эксплуатационные характеристики продукции.

\subsection*{Классификация задач выявления аномалий}
Согласно \cite{Hutten2024}, задачи выявления аномалий можно разделить на несколько категорий:
\begin{itemize}
	\item Обнаружение повреждений (defect detection). Эта задача охватывает все случаи, когда требуется классифицировать или выявить один или несколько типов дефектов. Особым подтипом здесь выделяют обнаружение трещин, что связано с высокой частотностью упоминаний этого случая в научной литературе. Данные методы применяются в диагностике строительных конструкций, мониторинге дорожных покрытий и тротуаров, выявлении дефектов бетона и асфальта. Формулировки задач включают классификацию, локализацию и сегментацию трещин.
	\item Проверка наличия или отсутствия деталей (completeness check). Эта категория задач направлена на определение отсутствующих элементов изделия или, наоборот, подтверждение их присутствия. Она охватывает широкий спектр применений, связанных с контролем сборочных операций и проверкой полноты комплектации.
\end{itemize}
Для решения подобных задач традиционно применяются классические методы машинного зрения: фильтрация, выделение признаков, пороговая обработка, морфологические операции и др. Однако в последние годы наблюдается стремительный рост числа работ, посвящённых использованию методов глубокого обучения. Во многих исследованиях применяется комбинированный подход, объединяющий глубокие нейронные сети и традиционные методы обработки изображений для повышения точности и надёжности решений.

\subsection*{Структура обзора}
\begin{enumerate}
	\item Классификация задач визуального контроля — первая секция посвящена систематизации разновидностей задач, связанных с обнаружением дефектов и аномалий.
	\item Традиционные методы обработки изображений — приводится описание классических подходов и их применимость к задаче выявления аномалий.
	\item Методы глубокого обучения — подробно рассматриваются современные подходы, основанные на нейронных сетях, их архитектуры, особенности применения и преимущества в сравнении с традиционными методами.
	\item Методы глубокого обучения, основанные на обучении без учителя.
\end{enumerate}
Over the years, defect detection methodologies have evolved from traditional manual inspection practices to sophisticated automated systems, driven by advancements in computer vision and deep learning.



\section{Традиционные методы машинного обучения}
Автоматические системы обнаружения дефектов были предложены в 1980-х годах \cite{Hutten2024}. Постепенно, с широким распространением визуальных сенсоров, появлялось всё больше методов контроля, основанных на машинном обучении (ML) и техниках обработки изображений, особенно в период с 1981 по 2010 годы. Такие методы извлекают заранее разработанные вручную признаки (handcraft features), которые затем подаются в распознаватель образов, поэтому их также называют методами feature engineering. 

Обычно традиционные методы на основе ML делятся на четыре типа: статистические методы, структурные методы, методы на основе фильтров и методы на основе моделей \cite{Hutten2024, Qifan2022}. Ниже приведено краткое описание этих типов согласно обзору \cite{Qifan2022}.

\subsection{Статистические методы}
Согласно \cite{Qifan2022}, статистические методы объединяют подходы, применяемые для анализа текстуры и выявления дефектов на основе статистических характеристик распределения значений пикселей в изображении. Вычисление текстурных признаков в этих методах выполняется покомпонентно (поточечно), что обеспечивает высокую устойчивость к изменениям освещения и воздействию шумов.

Статистические методы можно разделить на три основные категории:
\begin{enumerate}
	\item Методы статистики по уровням серого (grayscale statistics). Текстурные признаки извлекаются путём анализа значений полутонов пикселей. Наиболее распространённые показатели включают среднее значение, дисперсию, гистограммы, моменты распределения, энтропию и другие характеристики, отражающие структуру распределения интенсивностей как в целом, так и в локальных областях изображения.
	\item Методы статистики по шаблонам (pattern statistics). Эти методы описывают локальные структурные паттерны, основанные на взаимосвязях между пикселем и его ближайшими соседями. Часто они представляют собой формирование двоичных или категориальных кодов, характеризующих локальную структуру текстуры. Типичными примерами являются методы семейства Local Binary Patterns (LBP) и другие схожие подходы.
	\item Методы статистики по градиентам (gradient statistics). Основываются на анализе локальных изменений интенсивности, то есть градиентов по направлениям X и Y. Такие методы учитывают величину и ориентацию градиента, что делает их особенно эффективными для обнаружения краёв, линий, трещин и царапин.
\end{enumerate}
Несмотря на устойчивость к шумам и освещению, статистические методы имеют ограничение: они не учитывают глобальную структуру изображения и не позволяют полноценно исследовать взаимосвязи между удалёнными пикселями. Это снижает их эффективность при анализе сложных текстур и крупномасштабных аномалий.

\subsection{Спектральные методы (в других источниках - методы на основе фильтров)} 
Для некоторых видов дефектов, которые сложно обнаружить только по яркости пикселей (градациям серого), характерно нарушение структурной однородности равномерных текстур. Чтобы выявить такие нарушения, изображение можно преобразовать из временной области в частотную. В частотном представлении дефекты становятся лучше различимы благодаря заметным отличиям отклика нормальной и повреждённой текстуры.

К методам частотного анализа относятся:
\begin{enumerate}
	\item Преобразование Фурье (Fourier Transform). Позволяет перейти от временной области к частотной и разделить сигнал по частотам. Метод эффективен для стационарных и детерминированных сигналов, относится к глобальным подходам анализа и хорошо подходит для изучения стабильных структур на больших интервалах.
	\item Фильтр Габора (Gabor Filter). Выполняет локальный частотно-временной анализ с фиксированным разрешением. Оно особенно чувствительно к краям объектов и мало зависит от освещения, что делает его полезным для обнаружения локальных текстурных нарушений.
	\item Вейвлет-преобразование (Wavelet Transform). Обеспечивает гибкий и мощный локальный анализ сигналов. Оно позволяет менять форму окна и спектральную структуру, благодаря чему эффективно выявляет резкие изменения, особенности и сингулярности в изображении.
\end{enumerate}

\subsection{Структурные методы}
Структурные подходы (Structural Methods) сосредоточены на пространственном расположении элементов текстуры. Эти элементы можно извлечь из текстуры и описать как текстурные примитивы. Применяя правила пространственного расположения к этим примитивам, можно получить динамическую модель текстуры.

Как правило, текстурные примитивы представляют собой простые области серого цвета, отрезки линий или отдельные пиксели. Эти элементы всегда используются совместно с правилами размещения, которые выводятся из геометрических взаимосвязей или пространственной статистики этих примитивов.

Согласно обзору \cite{Czimmermann2020} структурные методы показывают значительно лучшие результаты при работе с регулярными (упорядоченными) текстурами.

\subsection{Методы на основе моделей}
Методы анализа текстур на основе моделей стремятся описать процесс генерации текстур. Их цель — смоделировать текстуру, определив параметры заранее заданных моделей. Такой подход учитывает не только локальную случайность текстуры, но и глобальную закономерность её структуры, что делает метод более гибким. Однако основными его недостатками являются сложность оценки параметров модели и высокая вычислительная трудоемкость.

Модельные методы подразделяются на три группы \cite{Czimmermann2020}:
\begin{enumerate}
	\item Фрактальные модели. Фракталы играют важную роль в описании естественных поверхностей с самоподобной и нерегулярной текстурой. Впервые они были представлены Мандельбротом. Фрактальные методы основаны на вычислении и анализе фрактальных чисел и других фрактальных характеристик. Широко распространенным является метод подсчета коробок (box-counting) и его модификации. Эффективность методов этой категории сильно зависит от степени самоподобия текстуры, и, поэтому при анализе не самоподобных структур, метод этой категории оказываются менее устойчивыми.
	\item Авторегрессионные модели. Основная концепция авторегрессионной модели (AR) заключается в характеристике признаков текстуры на основе линейных зависимостей пикселей \cite{Czimmermann2020}.
	\item Модели случайных полей. Подходы на основе случайных полей Маркова сочетают статистическую и структурную информацию контекстно-зависимых сущностей, таких как пиксели, зависящие от соседних пикселей, и могут использоваться в задачах сегментации и классификации текстур.
\end{enumerate}

\section{Классификация задач} %Классификация по типу области применения/Classification based on Application Requirements and Context

В обзоре \cite{Bhatt2021} приведена классифкация различных задачи обнаружения дефектов на основе прикладных требований и контекста. Авторы обзора делят задачи на
\begin{enumerate}
	\item Обнаружение аномалий
	\item Целевое обнаружение дефектов
	\item Одновременная идентификация множественных дефектов
	\item Кластеризация по типам дефектов. 
\end{enumerate}
Ниже приведено краткое изложение выделенных типов задач.

\subsection{Обнаружение аномалий}
Согласно исследованию \cite{Bhatt2021}, обнаружение аномалий при контроле дефектов представляет собой процесс идентификации отклонений в наборах данных, которые соответствуют дефектам. Данная задача обычно рассматривается в контексте неконтролируемого обучения [41], где аномалии определяются как непредвиденные события, отклоняющиеся от нормальных данных. В последние годы наблюдается растущий интерес к применению глубокого обучения для решения этой задачи.

Методы обнаружения аномалий на основе глубокого обучения можно классифицировать на три основные категории:

\subsubsection{Контролируемое обнаружение аномалий}
Данный подход использует размеченные обучающие выборки, содержащие как бездефектные, так и дефектные образцы. Хотя этот метод может демонстрировать высокие показатели обнаружения благодаря полной размеченности данных, его эффективность часто ограничивается проблемой дисбаланса классов в наборах данных.

\subsubsection{Полуконтролируемое обнаружение аномалий}
Также известный как одноклассовая классификация, этот метод использует только размеченные нормальные образцы для обучения. Основная идея заключается в определении дискриминационной границы, включающей бездефектные образцы, при этом любые образцы за пределами этой границы классифицируются как аномалии. Преимущество подхода - отсутствие необходимости сбора большого количества дефектных образцов. Однако точность метода может уступать контролируемому подходу. Исследования [43-46] демонстрируют различные реализации этого метода, причем работа [46] показывает более четкие границы принятия решений для выбросов.

\subsubsection{Неконтролируемое обнаружение аномалий}
Применяется когда размеченные данные недоступны и основывается на двух ключевых предположениях:
\begin{itemize}
	\item Большинство образцов в наборе данных являются нормальными
	\item Признаки аномальных образцов значительно отклоняются от стандартных
\end{itemize}

Метод изучает внутренние характеристики набора данных для отделения аномалий. В работе [47] демонстрируется применение генеративных состязательных сетей (GAN), обученных только на "здоровых" данных. Другой пример - метод Deep Support Vector Data Description (DeepSVDD) [48], который минимизирует объем гиперсферы, содержащей сетевые представления нормальных данных, классифицируя внешние точки как аномалии.

\subsection{Целевое обнаружение дефектов (Targeted Defect Detection)}

Метод целевого обнаружения дефектов позволяет выявлять конкретный тип дефекта, заданный в качестве целевого. Сначала идентифицируется целевой дефект, а затем определяется его местоположение на изображении. Этот метод реализуется с использованием обучения с учителем, поскольку конкретная цель должна быть помечена как дефект. Для обучения модели требуется значительное количество дефектных образцов. Сложность вычислений зависит от типа дефекта и объема собираемых данных [49]. 

В отличие от контролируемого обнаружения аномалий, которое выявляет отклонения в данных после обучения на хороших образцах, целевое обнаружение дефектов обучает модель или глубокую нейронную сеть на дефектных образцах. Процесс сбора данных может быть трудоемким из-за необходимости получения достаточного количества дефектных образцов. Например, в стандартных промышленных условиях и производственных линиях количество образцов без дефектов обычно значительно превосходит количество дефектных образцов. Поэтому сбор большого количества образцов конкретного дефекта может занять время, и такой сценарий может быть не идеальным. Однако этот метод хорошо работает, когда данные о дефектах легко доступны. Обнаружение трещин в дорожных покрытиях может служить репрезентативным примером [50], где CNN используется для детектирования трещин на асфальте.

\subsubsection{Особенности реализации}
\begin{itemize}
	\item Требуется обучение с учителем, поскольку целевой дефект должен быть предварительно размечен
	\item Для эффективного обучения необходим значительный объем дефектных образцов
	\item Сложность вычислений зависит от типа дефекта и объема данных [49]
\end{itemize}

\subsubsection{Сравнительный анализ}
В отличие от контролируемого обнаружения аномалий, которое обучается на нормальных образцах, целевое обнаружение ориентировано именно на дефектные данные. Это создает определенные проблемы:
\begin{itemize}
	\item Трудоемкость сбора данных из-за необходимости значительного количества дефектных образцов
	\item Дисбаланс в производственных условиях: нормальные образцы обычно преобладают
	\item Временные затраты на сбор репрезентативной выборки конкретного дефекта
\end{itemize}

Несмотря на эти ограничения, подход демонстрирует эффективность при доступности достаточного количества данных о дефектах.

\subsubsection{Практический пример}
Обнаружение трещин в дорожном покрытии [50] является характерным примером успешного применения метода с использованием сверточных нейронных сетей (CNN). \cite{Bhatt2021}

\subsection{Конкурентная идентификация множества дефектов (Concurrent Identification of Multiple Defects)}

Поверхностные дефекты могут иметь разнообразные формы, такие как царапины, трещины, включения, пятна, вмятины, отверстия и многие другие. Контролируемая поверхность может содержать один или несколько таких дефектов одновременно. Чтобы определить причину возникновения множественных дефектов, необходимо сначала идентифицировать их все. Можно применять целевое обнаружение дефектов для каждого из множественных дефектов по отдельности. Однако этот подход не будет эффективным, поскольку требует отдельно обученных моделей глубокого обучения для каждого дефекта. Кроме того, он может давать сбои, если существует перекрытие между дефектами. Например, на дефектной поверхности трещина может пересекать вмятину. В производственных условиях часто наблюдается одновременное присутствие нескольких типов дефектов, что требует разработки специализированных методов идентификации.

Подход для совместной идентификации множественных дефектов предполагает использование единой архитектуры глубокого обучения для определения типа дефекта на рассматриваемой поверхности. Это можно выполнить в один этап или разделить на два шага: 
\begin{enumerate}
	\item Определить, есть ли на поверхности дефект
	\item Классифицировать тип дефекта
\end{enumerate}
Множественные дефекты, идентифицируемые этим подходом, включают только известные дефекты. Каждый из них может быть размечен вручную или автоматически, но они должны быть помечены в соответствии с известными категориями поверхностных дефектов. Таким образом, такой подход возможен только с использованием обучения с учителем или полуконтролируемого обучения. Размеченные изображения используются для обучения модели глубокого обучения с целью распознавания известных дефектов. Обучение должно включать образцы поверхностей с множественными дефектами и их взаимодействиями, чтобы модель могла предсказывать дефекты совместно.

{\color{red} Проблема совместной идентификации множественных дефектов изучалась в [34]. Вся архитектура системы разделена на четыре этапа: (1) обнаружение аномалий, (2) фильтрация ложных аномалий, (3) кластеризация пикселей дефектов и (4) классификация дефектов. Она следует двухэтапной модели обнаружения аномалий и классификации дефектов с добавлением этапов фильтрации и кластеризации.}

Ограничение подхода совместной идентификации — необходимость в размеченных данных о дефектах для обучения. Некоторые дефекты встречаются редко, и по ним недостаточно данных или они не размечены. Это может происходить при адаптации нового производственного процесса. В таких случаях дефекты могут остаться необнаруженными этим методом. Однако поскольку эти методы могут обнаруживать дефекты совместно, они позволяют выявлять взаимодействия между дефектами. Это, в свою очередь, может позволить обнаружить взаимозависимость между множественными поверхностными дефектами. Существует необходимость в исследованиях области использования совместного обнаружения дефектов для выявления взаимосвязей между поверхностными дефектами.

\subsection{Кластеризация по типу дефекта (Defect-Type Clustering)}

В настоящее время работ по неконтролируемой классификации дефектов минимальное количество. Одним из возможных подходов является кластеризация по типам дефектов. Это улучшит подход к классификации поверхностных дефектов и позволит работать с неограниченным количеством типов дефектов. Это также устранит трудоемкий процесс маркировки типов дефектов для всех обучающих образцов. Исследователям следует изучить эту область для будущих исследований в области глубокого обучения для обнаружения поверхностных дефектов.

В реальных промышленных условиях существует значительное разнообразие сценариев, когда традиционные методы контроля дефектов оказываются неэффективными:
\begin{itemize}
	\item Наличие редких и каскадных дефектов
	\item Появление новых типов дефектов при изменении технологических процессов
	\item Ситуации с теоретически неограниченным разнообразием морфологических характеристик дефектов
\end{itemize}
Контролируемые и полуконтролируемые методы, описанные в обзоре \cite{Bhatt2021}., демонстрируют ограниченную эффективность в таких условиях, что создает потребность в разработке неконтролируемых методов кластеризации.

Предлагаемый подход включает следующие этапы:
\begin{enumerate}
	\item Первичное обнаружение аномалий без предварительной классификации.
	\item Автоматическая кластеризация дефектов по морфологическим признакам
	\item Извлечение характеристик (форма, размер, цвет, текстура)
	\item Интерактивная классификация с присвоением условных обозначений (тип 1, тип 2, и т.д.)
\end{enumerate}

Исследователи обсуждали схожий неконтролируемый подход в своем обзоре по визуальному контролю стальной продукции [51]. Однако эта работа ограничивается использованием самоорганизующихся карт (SOM) искусственной нейронной сети для классификации множественных дефектов на поверхности стали и не рассматривает методы визуального контроля на основе глубоких нейронных сетей. Другие исследователи применили неконтролируемый подход для классификации дефектов [52]. Они использовали два автоэнкодера и слой классификации с функцией активации Softmax в своей модели глубокого обучения. Автоэнкодеры всегда обучаются неконтролируемым способом, но они обучали слой Softmax контролируемым способом для классификации дефектов поверхности стали. Рассматриваемые ими различные типы дефектов были взяты из базы данных NEU по дефектам поверхности стали. Если бы они обучили слой Softmax с использованием неконтроливаемого подхода, это был бы подход к кластеризации дефектов.

\subsubsection*{Преимущества и ограничения подхода}
Преимущества:
\begin{itemize}
	\item Отсутствие необходимости в размеченных данных
	\item Возможность обработки неизвестных типов дефектов
	\item Автоматизация процесса классификации
\end{itemize}

Ограничения:
\begin{itemize}
	\item Недостаточная проработка глубоких нейросетевых подходов
	\item Ограниченность существующих методов (например, SOM)
	\item Сложности валидации результатов кластеризации
\end{itemize} 

Кластеризация по типам дефектов представляет значительный интерес для развития систем автоматизированного контроля качества. Несмотря на ограниченное количество текущих исследований, этот подход обладает потенциалом для решения сложных задач промышленного контроля, особенно в условиях неопределенности и постоянно меняющихся производственных процессов. Дальнейшие исследования в этом направлении могут существенно расширить возможности систем технического зрения в промышленных приложениях.

\section{Классификация методов на основе глубокого обучения (Learning-Based Classification)}
Глубокое обучение является одной из самых быстрорастущих областей в информатике благодаря своей способности решать чрезвычайно сложные задачи. Накопление богатого арсенала традиционных методов машинного обучения привело к эволюции глубокого обучения, которое также черпает вдохновение в статистическом обучении. Большинство подходов, упомянутых в предыдущих разделах, считаются традиционными решениями, где основное внимание уделяется явно сконструированным признакам, которые может быть сложно описать в сложных случаях. Однако глубокое обучение использует обучение представлениям данных для выполнения задач, преобразуя данные в сложные, абстрактные представления, которые позволяют системам обучаться признакам (например, обучение признакам). Эта способность глубокого обучения устраняет необходимость в сложных признаках для конкретного дефекта. Как глубокое обучение, так и традиционное машинное обучение являются ориентированными на данные методами искусственного интеллекта, способными успешно моделировать детерминированные правила, которые часто непостижимы для человека, и отношения между входными и выходными данными. Более того, глубокое обучение обладает возможностью выполнять обучение признакам, построение модели и обучение модели путем выбора различных ядер или настройки и оптимизации параметров \cite{Czimmermann2020}.

\subsubsection*{Преимущества перед традиционными методами}

Глубокое обучение — это быстроразвивающаяся область компьютерных наук, которая демонстрирует высокую эффективность в решении сложных задач. В отличие от традиционных методов машинного обучения, где особенности (features) должны быть сконструированы вручную (что особенно трудно для сложных дефектов), глубокое обучение автоматически извлекает и обучается этим особенностям из самих данных. Этот процесс, известный как обучение представлениям (representation learning), преобразует сырые данные в сложные абстрактные представления, что устраняет необходимость в сложном инженерном проектировании признаков для каждого конкретного типа дефекта \cite{Czimmermann2020}.

\subsubsection*{Ключевые проблемы и вызовы}

Несмотря на потенциал, внедрение глубокого обучения сталкивается с несколькими серьезными проблемами:

Нехватка размеченных данных: Создание больших наборов данных, специфичных для конкретной задачи, является дорогостоящим и трудоемким процессом. Проблема усугубляется тем, что дефекты на производстве часто являются редкими событиями.

Требование точной локализации: В отличие от простой классификации изображений, во многих промышленных задачах требуется не только определить наличие дефекта, но и точно очертить его границы (сегментация) для последующего ремонта или анализа.

Для преодоления этих вызовов исследователи разрабатывают и используют различные методики \cite{Bhatt2021}:
\begin{itemize}
	\item Аугментация данных (Data Augmentation): Искусственное расширение набора данных путем применения к изображениям различных преобразований (повороты, масштабирование, искажения).
	\item Трансферное обучение (Transfer Learning): Использование предварительно обученных моделей (например, на наборе данных ImageNet) и их дообучение на небольшом целевом наборе данных.
	\item Специализированные архитектуры CNN: Модификация и разработка новых структур нейронных сетей, лучше подходящих для задач обнаружения и точной сегментации дефектов.
	\item Комбинирование подходов: Интеграция контролируемых и неконтролируемых методов в рамках единого конвейера для повышения эффективности при недостатке данных и решения задачи точного очерчивания дефектов.
\end{itemize}

 В данном разделе более подробно описаны типы моделей глубокого обучения в применении к задаче выявления аномалий изображений. Классификация приведена из обзора \cite{Bhatt2021}.

\subsection{Обучение с учителем (Superwised Learning)} 
В задачах обучения с учителем (supervised learning) есть набор «правильных ответов», и нужно его продолжить на все пространство или, точнее говоря, на те примеры, которые могут встретиться в жизни, но которые еще не попадались в тренировочной выборке; к этому классу относятся задачи регрессии и классификации, как например классификация рукописных цифр по набору данных MNIST или машинный перевод на основе параллельных текстов на разных языках; \cite{Nikolenko2018}

Обучение с учителем предполагает подачу на вход алгоритму обучения пар входных данных \(\{\mathbf{x}_{i}\}\) и соответствующих им целевых выходных значений \(\{t_{i}\}\). Алгоритм настраивает параметры модели так, чтобы максимизировать согласие между прогнозами модели и целевыми выходными данными. Выходные данные могут быть либо дискретными метками из набора классов \(\{\mathcal{C}_{k}\}\), либо набором непрерывных, потенциально векторных величин, которые мы обозначаем как \(\mathbf{y}_{i}\), чтобы четче разграничить эти два случая. Первая задача называется классификацией поскольку мы пытаемся предсказать принадлежность к классу, тогда как вторая называется регрессией, поскольку исторически аппроксимация тренда по данным называлась именно так.

После фазы обучения, в течение которой все учебные данные (размеченные пары "вход-выход") обрабатываются (часто путем многократного их перебора, обученная модель может теперь использоваться для прогнозирования новых выходных значений для ранее не встречавшихся входных данных. Эта фаза часто называется тестовой фазой, хотя это иногда заставляет людей чрезмерно сосредотачиваться на производительности на конкретном тестовом наборе, вместо того чтобы создавать систему, которая надежно работает для любых правдоподобных входных данных, которые могут возникнуть. \cite{Szeliski2022}

Достоинства подхода:
\begin{itemize}
	\item Высокая точность распознавания дефектов
	\item Возможность классификации дефектов
\end{itemize}

Недостатки подхода:
\begin{itemize}
	\item Высокая вычислительная сложность
	\item Зависимость от больших объемов размеченных данных
	\item Высокая стоимость и сложность разметки данных
	\item Не способны адаптироваться к новым типам данных \cite{Li2025} Невозможность обнаружить новые, неизвестные типы дефектов.
	\item Дисбаланс классов
	\item Сложность адаптации к изменениям
	\item Риск переобучения
\end{itemize}

\subsection{Полуавтоматическое обучение (Semi-supervised Learning)}

In many machine learning settings, we have a modest amount of accurately labeled data and a far larger set of unlabeled or less accurate data. For example, an image classification dataset such as ImageNet may only contain one million labeled images, but the total number of images that can be found on the web is orders of magnitudes larger. 

Semi-supervised learning is a subset of the larger class of weakly supervised learning problems, where the training data may not only be missing labels, but also have labels of questionable accuracy (Zhou 2018). Some early examples from computer vision (Torresani 2014) include building whole image classifiers from image labels found on the internet (Fergus, Perona, and Zisserman 2004; Fergus, Weiss, and Torralba 2009) and object detection and/or segmentation (localization) with missing or very rough delineations in the training data (Nguyen, Torresani et al. 2009; Deselaers, Alexe, and Ferrari 2012). In the deep learning era, weakly supervised learning continues to be widely used (Pathak, Krahenbuhl, and Darrell 2015; Bilen and Vedaldi 2016; Arandjelovic, Gronat et al. 2016; Khoreva, Benenson et al. 2017; Novotny, Larlus, and Vedaldi 2017; Zhai, Oliver et al. 2019). A recent example of weakly supervised learning being applied to billions of noisily labeled images is pre-training deep neural networks on Instagram images with hashtags (Mahajan, Girshick et al. 2018).\cite{Szeliski2022}

Полуконтролируемые методы обнаружения аномалий используют небольшое количество размеченных данных совместно с большим объёмом неразмеченных изображений. Это позволяет повысить точность обнаружения при минимальных затратах на разметку \cite{Shukla2025}. В таких подходах применяются механизмы самотренировки, регуляризация устойчивости, графовые методы и генеративные модели (GAN, автоэнкодеры, диффузионные модели), что делает их эффективными при дефиците разметки (Saheel et al., 2024; Saeedi and Giusti, 2023).

\subsubsection*{Примеры полуконтролируемых методов}
Di [56] предложил подход на основе сверточного автоэнкодера и полу-контролируемых GAN. Сначала на неразмеченных данных обучается CAE, затем его энкодер используется в составе GAN-дискриминатора и дообучается на размеченных образцах.

He [57] применил схему, в которой GAN генерирует новые примеры, а остаточная сеть присваивает им псевдометки; затем модель многократно переобучается на расширенном наборе данных. Метод достигает высокой точности даже при очень малом количестве исходных образцов.

Gao [58] использовал CNN с методом Pseudo-Labeling, который автоматически присваивает метки неразмеченным данным.

Xie [59] реализовал полу-контролируемую DCGAN, где общие сверточные слои используются одновременно дискриминатором и классификатором. Потеря комбинирует кросс-энтропию и расстояние Вассерштейна, что повышает устойчивость при обучении.

Zhang [60] предложил SSGAN с двумя подсетями: сегментационной сетью с двойным механизмом внимания и полносверточным дискриминатором, который генерирует карты доверия для неразмеченных изображений. Метод улучшает сегментацию и снижает зависимость от разметки.

\subsubsection{Стратегии преодоления нехватки данных}
Полуконтролируемые и неконтролируемые методы часто используют аугментацию, перенос обучения и предварительную обработку.
В [69] увеличивают число примеров, вырезая ключевые фрагменты дефектов.

В [70] и [55] применяют transfer learning на базе AlexNet и Decaf, заменяя финальный классификатор на двухклассовый.

В [71] используется SqueezeNet, в [72] — автоэнкодер для сравнения дескрипторов дефектов без разметки.

В [73] обучение ведётся только на бездефектных образцах — формируется эталонное представление, с которым затем сравниваются тестовые изображения.

Когда число классов дефектов велико или возникают новые типы повреждений, ключевым инструментом становится обнаружение аномалий, основанное на сравнении признаков тестового и эталонного образцов. В [58] также предлагается автоматизированная разметка данных на основе выделения областей дефектов.


\subsection{Обучение без учителя/Unsupervised Learning}

Самообучение (Self-supervised learning) — это парадигма машинного обучения, в которой модель учится прогнозировать определенные 

В задачах обучения без учителя (unsupervised learning) есть просто набор примеров без дополнительной информации, и задача состоит в том, чтобы понять его структуру, выделить признаки, хорошо ее описывающие. \cite{Nikolenko2018}

Self-supervised learning is a machine learning paradigm where a model learns to predict certain properties or features of input data without the need for explicit labels or annotations. Instead of relying on manually labeled data, the model exploits the inherent structure and patterns within the data to generate supervised like signals for training. The key idea is to make use of abundant unlabeled data to learn meaningful representations that can be useful in different tasks. 

The training process in self-supervised learning revolves around pretext tasks — artificial tasks that generate learning signals from the data itself. For instance, in computer vision, a model can be trained to predict the rotation angle of an image or to restore missing patches of an image. Such pretext tasks help the model develop an understanding of the data, which can then be leveraged for downstream tasks like classification or object detection. After training on the pretext task, the learned representations are fine-tuned on tasks requiring labeled data, making the model highly efficient and versatile.

In recent years, self-supervised learning for computer vision has gained significant popularity due to its ability to utilize large quantities of unlabeled data to enhance the performance of downstream tasks. This approach has been successfully implemented across various domains, including natural language processing, computer vision, and speech recognition, bridging the gap between supervised and unsupervised learning and advancing the capabilities of artificial intelligence in tackling real-world problems. \cite{MasteringCompVis}

Неконтролируемые методы обнаружения визуальных аномалий получают всё большее распространение, поскольку они не требуют размеченных данных и строят модели, используя только нормальные изображения (Lee and Kang, 2022; Sun, 2024; Tan and Wong, 2024). Исследования показывают, что неконтролируемые методы позволяют машине самостоятельно выявлять внутренние структуры данных и классифицировать изображения без разметки. Это делает их удобными для промышленности, где получить дефектные образцы сложно и дорого. Такие методы выявляют широкий спектр отклонений, включая новые и неожиданные дефекты, что повышает надежность систем контроля (Pinon et al., 2024; Zhang F., 2023).

Основные подходы включают:
\begin{enumerate}
	\item Методы реконструкции. Основаны на том, что модели, обученные на нормальных данных (AEs, VAEs, GANs), плохо восстанавливают аномальные области. Ошибка реконструкции служит метрикой для их выявления (Lin, 2023).
	\item Методы на основе нормализующих потоков. Используют обратимые нейросети для моделирования распределения нормальных данных. Аномалии определяются по низкой вероятности в обученном распределении.
	\item Методы представлений (representation-based). Изучают устойчивые признаки нормальных объектов. Отклонение от таких представлений трактуется как аномалия. Используются deep metric learning и self-supervised learning (Zingman et al., 2024).
	\item  Методы с аугментацией данных. Усиливают устойчивость моделей путем добавления различных преобразований нормальных изображений, генерации синтетических данных и adversarial-тренировок (Oyelade and Ezugwu, 2021; Motamed et al., 2021).
\end{enumerate}

\section{Defect Inspection Based on Generative Models}
Поскольку исходная задача относится к категории задач, решаемых с помощью самообучающихся моделей, далее приведено описание наиболее распространенных типов моделей такого типа.

При проверке дефектов современные методы получают сгенерированные образцы с помощью четырех типов генеративных моделей: VAEs, GANs, DMs и мультимоделей, которые комбинируют несколько моделей для генерации образцов. Кроме того, задачи инспекции дефектов можно разделить на три основные задачи компьютерного зрения:  
\begin{enumerate}
	\item Классификация дефектов — определение класса дефекта для образца;
	\item Обнаружение дефектов — классификация и локализация каждого экземпляра дефекта в образце;
	\item Сегментация дефектов — разделение образца на области с однородными характеристиками пиксель за пикселем.
\end{enumerate}

В этом разделе представлен технический обзор методов инспекции дефектов на основе генеративных моделей, предложенных в период с 2022 по 2024 год.\cite{YuHe2024}

\subsection{Autoencoders and Variational Autoencoders}

Вариационные автоэнкодеры (VAE) широко применяются в инспекции поверхностных дефектов [32]. Наиболее простой способ их использования — генерация (augmentation) дополнительных данных, которые затем подаются в классификатор. Например: Сян и др. [56] использовали сверточный автоэнкодер для генерации дефектов тканей и применили метод обнаружения на основе преобразования Фурье, Чжан и др. [57] разработали AE5-SSIM, реконструирующий изображения фольги с учётом структурного сходства, Сае-анг и др. [58] применили автоэнкодер для извлечения признаков из немаркированных данных и последующего семантического анализа, Си и др. [59] использовали VAE для реконструкции изображений деревянных поверхностей и исследовали влияние различных фильтров градиента на качество обнаружения, Ши и др. [60] предложили облегчённую CAE-сеть для решения задачи нехватки дефектных образцов.

Однако подход на основе генерации данных имеет ограничения: раздельное обучение моделей генерации и классификации не позволяет учитывать взаимосвязи признаков. Поэтому современные исследования стремятся к архитектурам со сквозным (end-to-end) обучением.

К таким разработкам относятся:
\begin{itemize}
	\item AVQ-VAE с механизмом внимания [61];
	\item много масштабная U-образная модель MSUDCAE, использующая карты невязок между оригиналом и реконструкцией [62];
	\item неконтролируемая двойная Siamese-сеть для инспекции аномалий [63];
	\item двухэтапная архитектура DR-VAE для анализа железнодорожных поверхностей [64];
	\item VAE с улучшенными семантическими метками для генерации дефектов и повышения точности классификации [65];
	\item FastRecon — sample-less метод обнаружения аномалий в условиях дефицита данных [66];
	\item ISRM — самообучающийся интроспективный автоэнкодер [67];
	\item Fs-CAE, использующий частотный сдвиг для работы с бездефектными данными [68];
	\item ViT-архитектура для обнаружения дефектов на основе глобального контекстного анализа [69].
\end{itemize}

\subsubsection*{Ограничения VAE}
Несмотря на широкое применение, методы на основе VAE имеют существенные ограничения. Они наиболее эффективны в простых случаях, когда каждое изображение содержит только один дефект. Процесс кодирования-декодирования в VAE предполагает гауссово распределение входных данных и их отображение в непрерывное пространство [28], что хорошо работает для регулярных структур, но недостаточно для сложных и случайных признаков. Это приводит к проблеме генерации размытых изображений [32].

В сравнении с VAE, генеративно-состязательные сети (GAN) и диффузионные модели (DMs) обладают более сложными априорными предположениями о данных в латентном пространстве и демонстрируют лучшую способность работы со сложными наборами данных по дефектам [31]. Это делает их перспективным направлением для дальнейших исследований в области обнаружения поверхностных дефектов.

\subsubsection*{Теоретические основы AE и VAE}

Основная цель автоэнкодера — получение компактного и информативного представления данных. Модель состоит из кодировщика, формирующего скрытое представление (z(x)), и декодировщика, восстанавливающего выход (y(z)). Чтобы сеть не повторяла входные данные напрямую, вводятся ограничения: уменьшение размерности скрытого слоя, разреженность, либо необходимость восстанавливать «искажённые» данные (denoising).

Вариационный автоэнкодер (VAE) использует приближение правдоподобия через нижнюю оценку ELBO, применяет амортизированную оценку скрытых переменных и использует приём репараметризации для обеспечения дифференцируемости процесса обучения \cite{kingma2013autoencoding}.

\subsection{Generative Adversarial Networks}

Генеративно-состязательные сети (GAN) являются одним из наиболее распространённых классов генеративных моделей, применяемых для задач обнаружения и анализа дефектов промышленной продукции. В рамках существующих исследований выделяют две основные архитектурные стратегии: классические GAN-модели и методы на основе image-to-image трансформации.

\subsubsection{Классические архитектуры GAN}

Генеративно-состязательные сети (GAN) занимают ведущее положение среди генеративных моделей, применяемых для задач контроля дефектов. В рамках данного подхода можно выделить две основные архитектурные парадигмы. Первая основана на классических GAN, осуществляющих безусловную генерацию образцов из распределения случайного шума. Наиболее репрезентативными моделями этого направления являются стандартный GAN [42], глубокий сверточный GAN (DCGAN) [70] и GAN с метрикой Вассерштейна (WGAN) [71].

Основные направления применения классических GAN включают:
\begin{itemize}
	\item Аугментация данных - генерация дополнительных образцов для расширения обучающих выборок
	\item Модификация дискриминаторов - адаптация выходных слоев под конкретные задачи контроля
	\item Полуконтролируемое обучение - генерация немаркированных образцов с последующим присвоением меток
\end{itemize}

\subsubsection*{Практические реализации и модификации}
Среди наиболее значимых разработок, основанных на классической архитектуре GAN для контроля дефектов, можно выделить:

Специализированные архитектуры:
\begin{itemize}
	\item Специализированная GAN Хэ и др. [72] для генерации немаркированных образцов
	\item Генеративно-состязательные сети с аномальными латентными переменными (ALGAN) Мурасе и др. [74]
	\item DefectGAN Ли и др. [76] с полностью сверточными сетями и марковскими дискриминаторами
\end{itemize}

Интеграция с механизмами внимания:
\begin{itemize}
	\item WGAN с механизмами внимания для распознавания дефектов горячекатаной стали [77]
	\item ADRGAN Ло и др. [83] с механизмами внимания по каналам и пространству
	\item Адаптивная GAN (AGAN) Ма и др. [84] с самовнимающимся энкодером
\end{itemize}

Оптимизационные подходы:
\begin{itemize}
	\item EDCGAN Цзинь и др. [79] с функциями активации ELU
	\item GAN с двойными дискриминаторами и регуляризацией Лю и др. [80]
	\item Модифицированный ACGAN Ван и др. [87] с расстоянием Вассерштейна
\end{itemize}	

\subsubsection{Архитектуры image-to-image трансляции}

Второе направление представляет GAN на основе image-to-image трансляции, генерирующие образцы на основе входных изображений. К этому классу относятся CycleGAN [91], StyleGAN [92] и SinGAN [93]. Данный подход не только решает проблему ограниченности данных, но и направляет процесс обучения за счет ограничений, накладываемых входными образцами.

Ключевые применения включают:
\begin{itemize}
	\item Генерация бездефектных шаблонов с использованием CycleGAN [94]
	\item Слабо контролируемая сегментация с модулями внимания [95]
	\item Качественный синтез дефектов с улучшенной StyleGAN [97]
	\item Структурное обнаружение дефектов на основе pix2pix и transfer learning [98]
\end{itemize}

\subsubsection{Перспективные разработки}
Современные исследования демонстрируют тенденцию к созданию комплексных решений:
\begin{itemize}
	\item Прогрессивная GAN (Pro-GAN) с механизмом самовнимания [100]
	\item Интеграция U-Net с pix2pix для генерации карт локализации [101]
	\item Production-and-Elimination GAN (PE-GAN) для имитации формирования дефектов [103]
\end{itemize}
Представленные подходы демонстрируют эффективность GAN в решении практических задач промышленного контроля, обеспечивая генерацию качественных синтетических данных и улучшая показатели обнаружения дефектов в условиях ограниченности реальных образцов.

\subsection{Diffusion Machines}

Хотя доказано, что качество генерации диффузионных моделей (DMs) превосходит таковое у GAN и VAE, производственная сфера, судя по всему, уделяет мало внимания использованию DM для контроля дефектов и генерации образцов. Танг и др. [104] представили метод неконтролируемого обучения для обнаружения дефектов на поверхностях тканей, который улучшает качество генерации изображений за счет оптимизации временного шага и использует диффузионную модель с шумом Симплекса (SN-DDPM) для реконструкции изображений дефектов, реализуя тем самым их точную сегментацию. Хуанг и др. [105] предложили гибридный framework глубокого обучения на основе DDPM для augmentation данных и глубокой остаточной нейронной сети (DenseNet) для классификации дефектов на поверхностях композитных панелей. Тао и др. [106] генерировали разнообразные дефектные изображения, обучая модель DDPM на ограниченном наборе образцов для обобщенного контроля поверхностных дефектов. В этом методе предложен механизм стирания-восстановления (erasing–inpainting) для стимулирования рекомбинации реальных и изученных признаков, что позволяет полностью использовать ограниченные выборки дефектов. Лян и др. [107] представили метод стабилизированного диффузионного усиления на основе улучшенной модели YOLOv5 для обнаружения дефектов в металлических штампованных деталях. Этот метод внедряет облегченную структуру bottleneck transformer и механизм много головного самовнимания (multi-head self-attention) для повышения скорости и точности обнаружения.

Несмотря на то, что DM обладают высокой производительностью и могут генерировать реалистичные изображения, сетевая структура DM является слишком громоздкой, а время семплирования слишком велико, что требует значительных вычислительных ресурсов и времени обучения. Существует возможность для снижения сложности обучения и модели DM, что в будущем может повысить их применимость в промышленных приложениях. \cite{YuHe2024}

\section{Наборы данных}


\section{Сравнение методов/Анализ}

Традиционные методы обнаружения поверхностных дефектов, основанные на глубоком обучении, требуют большого количества данных для обучения признаков. Однако в реальных промышленных условиях получить данные о дефектах затруднительно, а разметка таких данных требует значительных трудовых и материальных ресурсов.
  

Обучение без учителя используется для задач оценки плотности распределения, снижения размерности и кластеризации. Тем не менее, во многих случаях методы обучения без учителя демонстрируют более низкую эффективность по сравнению с методами обучения с учителем. \cite{Czimmermann2020}

Обнаружение дефектов поверхности промышленных изделий на основе традиционной обработки изображений. Для обучения модели не требуется большого количества размеченных данных, и она имеет низкую временную сложность обучения. Однако данный метод предъявляет высокие требования к условиям съемки изображений и обладает слабой адаптивностью.

Обнаружение дефектов поверхности промышленных изделий на основе глубокого обучения. В данной статье в основном рассматриваются методы обнаружения дефектов на основе глубокого обучения, применимые при малом количестве размеченных данных. Эти методы обладают высокой адаптивностью, но характеризуются значительной временной сложностью и требуют использования большого количества высокопроизводительных вычислительных модулей.

Традиционные методы. Преимущества: Низкие требования к данным о дефектах, высокая скорость обнаружения, меньшее время обработки. Недостатки: Низкая степень интеллектуальности, узкая область применения, сильная зависимость от условий inspection. Не подходят для обнаружения нерегулярных дефектов.

Глубокое обучение. Преимущества: Автоматическое извлечение признаков, эффективное представление характеристик, широкая область применения. Недостатки: Требует значительных ресурсов, включая большие наборы данных, высокие затраты на ручную разметку и большое количество вычислительных мощностей. \cite{Xiangjie2022}	

Традиционные методы обработки изображений эффективны для решения определенного класса задач. Однако эти методы плохо справляются с шумом, изменениями условий освещения и фонами со сложной текстурой. В последнее время глубокое обучение активно исследуется для автоматизации обнаружения дефектов.

\section{Проблемы и нерешенные вопросы}

Какие задачи для дальнейших исследований выделили авторы обзора \cite{Hutten2024}:
\begin{itemize}
	\item В какой степени методы, применяемые к 2D-изображениям, могут быть перенесены на 3D-изображения, содержащие информацию о глубине, или на другие типы данных — такие как облака точек?
	\item изучение временного разрыва между академическими исследованиями и промышленными применениями, который, предположительно, характерен не только для AVI. Причины этого, вероятно, многообразны, но можно ли сократить выявленный разрыв в два–три года? И какие шаги для этого необходимы?
	\item Ожидается, что Vision Transformers станут предпочтительными моделями для задач AVI. Однако большинство публикаций сосредоточено на тонкой настройке предобученных моделей в рамках контролируемого обучения, игнорируя их мощный потенциал самообучающихся (self-supervised) методов. Эти возможности уже продемонстрированы на эталонных наборах данных как для специфических задач (например, детекции объектов), так и для обучения более универсальных представлений, применимых к множеству задач. Перенос методов самообучения на задачи AVI — это перспективное направление исследований, заслуживающее повышенного внимания, особенно с учётом того, что доступные промышленные датасеты значительно меньше эталонных исследовательских наборов данных.
\end{itemize}

\subsection{Проблема малых выборок и различий текстур (Small sample problem and texture difference challenge)} 
\cite{Shukla2025}
Ограниченный объем выборок и выраженные вариации текстур представляют собой серьезные препятствия для точного обнаружения аномалий, поскольку модели, обученные на недостаточном количестве данных, нередко демонстрируют слабую способность к обобщению.
\subsubsection{Комбинация обучения распределению данных и методов аугментации (Combination of data distribution learning and data augmentation recommendation)}

Применение современных методов обучения распределению данных в сочетании с техниками аугментации может существенно смягчить указанные проблемы. Создание синтетически расширенных наборов данных — посредством вращений, отражений, изменений цветовой гаммы и других преобразований — позволяет расширить разнообразие условий, с которыми сталкивается модель. Это повышает устойчивость алгоритма и способствует росту точности обнаружения дефектов. Ряд исследований (Fullington et al., 2024; Xu M. et al., 2023) подтверждает эффективность совместного использования этих подходов.

\subsubsection{Рекомендации по использованию трансферного обучения(Transfer learning recommendation)}
Трансферное обучение предполагает использование моделей, предварительно обученных на смежных задачах, с последующей их донастройкой для решения задач обнаружения аномалий. Такой подход позволяет преодолеть проблему малых выборок за счёт использования знаний, полученных на больших объёмах данных, что приводит к повышению качества распознавания на небольших специализированных выборках.

Ограниченный объём данных и высокая вариативность текстур значительно усложняют задачу построения обобщающих моделей: глубокие нейронные сети, обученные на малых выборках, как правило, склонны к переобучению, что приводит к ухудшению практических результатов.

Одним из наиболее эффективных способов решения этой проблемы является трансферное обучение. Использование предварительно обученных моделей позволяет извлекать информативные признаки даже в условиях ограниченной разметки. Так, Chen et al. (2023b) показали эффективность адаптивного трансферного обучения при мониторинге многорежимных технологических процессов и обнаружении аномалий на паровых турбинах, улучшив способность моделей к обобщению при изменении условий эксплуатации.

Аналогичным образом, Liu et al. (2023d) предложили SimpleNet — эффективную архитектуру для обнаружения и локализации аномалий на изображениях, использующую трансферное обучение для снижения влияния дефицита обучающих данных. Maray et al. (2023) подтвердили применимость данного подхода в медицине, существенно повысив точность обнаружения падений при малом количестве доступных примеров.

Обзор, выполненный Iman et al. (2023), систематизировал современные методы глубокого трансферного обучения и выделил перспективные направления повышения адаптивности моделей к новым доменам. Zhao et al. (2024) расширили применение концепции на промышленную сферу, продемонстрировав эффективность методов контрастивного и трансферного обучения при контроле малых компонентов на сборочных линиях.

В производстве трансферное обучение также показало высокую применимость. Например, Shin (2024) предложил материал-адаптивную архитектуру для обнаружения аномалий в технологии дуговой аддитивной наплавки, где использование конкатенированных свойств материала улучшило качество обнаружения дефектов. Cheng et al. (2022) применили трансферное обучение для повышения точности сегментации дефектов на тканях с помощью модифицированной архитектуры UNet, что позволило добиться хороших результатов даже при ограниченной разметке.

Значительные успехи отмечены и в медицине: Lanjewar et al. (2024) объединили модели трансферного обучения с LSTM для более точной диагностики рака груди по ультразвуковым изображениям, а Ani et al. (2024) продемонстрировали эффективность такого подхода в задачах мультиклассовой классификации аномалий молочной железы.

Таким образом, применение трансферного обучения обеспечивает действенный способ преодоления ограничений, вызванных малым объёмом выборки и вариациями текстур. Использование заранее обученных моделей позволяет существенно улучшить способность к обобщению, снизить потребность в больших объёмах размеченных данных и обеспечить высокую точность обнаружения аномалий в различных областях. \cite{Shukla2025}


Transfer Learning позволяет извлекать общие знания из исходной задачи и применять их к изучению целевой задачи, тем самым повышая обобщающую способность модели и сокращая временные и ресурсные затраты, вызванные необходимостью использования большого количества размеченных данных для целевой задачи. Ren [48] предложил универсальный метод автоматического обнаружения поверхностных дефектов, требующий малого количества обучающих данных. Метод строит классификатор на основе признаков фрагментов изображений, где признаки переносятся из предварительно обученной сети глубокого обучения, а затем получает поэлементные прогнозы путем свертки обученного классификатора с входным изображением. Ferguson [49] предложил систему распознавания дефектов литья в рентгеновских изображениях на основе структуры Mask R-CNN. Система обнаружения дефектов одновременно выполняет обнаружение и сегментацию дефектов на входном изображении и подходит для различных задач обнаружения дефектов. Модель использует перенос обучения для снижения требований к обучающим данным и повышения точности прогнозирования обученной модели. Если конкретнее, модель сначала обучается с использованием двух больших наборов данных общедоступных изображений, а затем обученная сеть переносится на относительно небольшой набор данных рентгеновских снимков металлического литья. Gong [50] предложил модель обнаружения объектов на основе переноса обучения с использованием Domain Adaptive Faster R-CNN (DA Faster) для случая, когда размер дефектов на рентгеновских изображениях композитной структуры космического аппарата очень мал, а количество обучающих образцов невелико. На базе DA Faster в подмодуль извлечения признаков добавлена Feature Pyramid Network (FPN) для много масштабной адаптации признаков, а также используются стратегия малых якорей (small anchor) и ROI Align для помощи в локализации дефектов малого размера. HUANG [51] внедрил метод метрического обучения из обучения с малым числом примеров в область обнаружения дефектов и предложил метод метрического переноса обучения с малым числом примеров для решения проблемы необходимости большого количества обучающих выборок в методах глубокого обучения. Эксперименты показывают, что метрическое обучение с малым числом примеров позволяет сети достигать лучших результатов на меньшем количестве образцов путем обучения на несвязанных крупномасштабных наборах данных. LIU [52] предложил стратегию повторного использования знаний для обучения CNN-модели с целью повышения точности и устойчивости обнаружения дефектов. Путем внедрения переноса обучения на основе моделей и расширения данных, знания из других задач компьютерного зрения переносятся в задачу промышленного обнаружения дефектов для достижения высокой точности при ограниченном количестве обучающих образцов. \cite{Qifan2022}

\subsubsection{Optimize network structure recommendation}

Оптимизация архитектуры нейронной сети с учётом специфики задачи обнаружения аномалий является ключевым фактором повышения эффективности алгоритма. Грамотно спроектированная модель способна точнее учитывать особенности данных малых выборок и различия текстур, обеспечивая устойчивость работы в различных условиях. Одним из подходов к подобной оптимизации является автоматизированный поиск архитектур (Neural Architecture Search, NAS), который позволяет подбирать оптимальные структуры нейросети путём перебора различных вариантов и выбора наиболее эффективного решения для конкретной задачи обнаружения аномалий. В качестве альтернативы возможно ручное проектирование лёгких, но результативных архитектур, адаптированных под ограничения конкретного набора данных.

Cao Y. и др. (2023) предложили архитектуру Collaborative Discrepancy Optimization для надёжной локализации аномалий на изображениях, показав, что оптимизация структуры сети позволяет повысить точность и устойчивость обнаружения. В их работе подчёркивается важность тонкой настройки архитектуры для повышения чувствительности к признакам аномалии при сохранении приемлемых вычислительных затрат.

Wan и соавт. (2024) развили данный подход, интегрировав метод сингулярного спектрального анализа (Singular Spectrum Analysis, SSA) с оптимизированной моделью ResNet50-BiGRU для обнаружения и прогнозирования аномалий на изображениях. Их результаты демонстрируют, что сочетание сверточных сетей (ResNet50) с рекуррентными архитектурами (BiGRU) позволяет эффективно выявлять как пространственные, так и временные паттерны аномалий, делая процесс детектирования более точным и адаптивным.

Таким образом, применение методов оптимизации архитектур — будь то NAS или целенаправленная ручная настройка — позволяет системам обнаружения аномалий достигать более высокой точности при сохранении компактности моделей. Это особенно важно для приложений реального времени и систем с ограниченными вычислительными ресурсами, где одинаково критичны как эффективность, так и надёжность детектирования.

\subsection{Проблема обнаружения малых объектов}

Обнаружение малых аномалий или объектов на изображениях представляет собой сложную задачу из-за их минимального занимаемого пространства в пикселях, что затрудняет их отличиение от шума. Низкое разрешение, фоновый шум и вариации масштаба создают значительные сложности для обнаружения малых объектов. Маленькие объекты часто не имеют distinguishable характеристик, что затрудняет их differentiation от окружающего фона для моделей глубокого обучения.

\subsubsection{Рекомендации}
Для решения этой проблемы повышение разрешения изображений, использование методов multi-scale извлечения признаков и внедрение механизмов внимания могут значительно улучшить точность обнаружения.

Многоscale извлечение признаков позволяет моделям захватывать мелкие детали на разных разрешениях. Feature Pyramid Networks (FPN) и YOLO (You Only Look Once) широко используемые архитектуры, которые улучшают обнаружение малых объектов за счет использования hierarchical feature maps. Дополнительно, механизмы внимания улучшают фокусировку на релевантных признаках, подавляя фоновый шум, тем самым повышая надежность обнаружения. Chen et al. (2021) представили MAMA-Net (Multi-Scale Attention Memory Autoencoder Network) для обнаружения аномалий, показав, что комбинация внимания с multi-scale memory networks улучшает распознавание малых объектов в медицинской визуализации. Их подход подчеркивает преимущества multi-scale обучения для захвата subtle аномалий.

Xiang et al. (2023) расширили эту концепцию, предложив Multi-Scale Attention and Dilation Network для обнаружения малых дефектов. Интегрируя dilated свертки с механизмами внимания, их модель эффективно извлекает fine-grained признаки, улучшая локализацию дефектов на промышленных поверхностях.

Xiao et al. (2022) разработали multi-scale CNN с усиленным слиянием признаков и механизмом внимания для распознавания внешнего вида точечных сварных швов. Их работа показала, что слияние multi-resolution признаков позволяет сети лучше различать нормальные и дефектные сварные швы.

Yang et al. (2020) дополнительно улучшили неконтролируемую локализацию аномалий, incorporаting multi-scale memory modules в автоэнкодеры, усиливая способность модели обнаруживать малые отклонения в структурированных средах.

Совсем недавно Tao et al. (2024) представили метод для обучения multi-resolution признаков для неконтролируемой локализации аномалий на промышленных текстурированных поверхностях. Их подход использует multi-scale представления для обнаружения subtle различий в текстуре, значительно улучшая производительность в реальных производственных приложениях.

Интегрируя multi-scale извлечение признаков, механизмы внимания и методы повышения разрешения, модели обнаружения малых объектов могут достичь более высокой точности, делая их применимыми в таких областях, как медицинская визуализация, промышленный контроль и дистанционное зондирование. Эти достижения гарантируют, что даже самые малые аномалии или дефекты точно идентифицируются, улучшая общую надежность системы.

\subsection{Проблемы на уровне данных}

Несбалансированные наборы данных, где нормальные образцы значительно превосходят по количеству аномальные, могут смещать модель в сторону majority класса, снижая эффективность обнаружения аномалий. Разнообразные наборы данных использовались для изучения проблем, присущих промышленному обнаружению аномалий.

\subsubsection{Рекомендации}
Решение этой проблемы на уровне данных может включать такие методы, как oversampling minority класса (аномалий) или undersampling majority класса (нормальных образцов). Дополнительно, методы синтетической генерации данных, такие как использование GAN (Generative Adversarial Networks) для создания реалистичных аномальных образцов, могут помочь сбалансировать набор данных. Внедрение этих стратегий гарантирует, что модель получает достаточное обучение на аномалиях, улучшая ее способности обнаружения.

Zeiser et al. (2023) предоставили детальную оценку методов глубокого неконтролируемого обнаружения аномалий, подчеркивая data-centric подход для онлайн-инспекции. Их работа выделяет, как характеристики набора данных значительно влияют на производительность модели, усиливая потребность в эффективных методах балансировки данных.

Yang Y. et al. (2023) предложили подход к обнаружению аномалий на основе глубокого обучения, который извлекает репрезентативные латентные признаки для обработки случаев, когда аномалии являются low-discriminative или недостаточны по количеству. Их метод улучшает обнаружение аномалий за счет усиления извлечения признаков из ограниченных аномальных данных.

Mun et al. (2024) решали проблему дисбаланса данных в AI-driven Clinical Decision Support Systems (AI-CDSS), представив UAnoGAN, GAN-based фреймворк для обнаружения аномалий. Их исследование показало, что синтетически сгенерированные аномалии могут значительно повысить robustность модели, особенно в медицинских приложениях, где реальные аномальные образцы scarce.

Liu D. et al. (2023) представили Deep Attention SMOTE, метод аугментации данных, который использует learnable interpolation factor для генерации синтетических образцов для несбалансированного обнаружения аномалий в газовых турбинах. Этот подход улучшает возможности обнаружения, создавая разнообразные обучающие образцы, которые предотвращают смещение модели в сторону majority класса.

Помимо простых методов resampling, Cao T. et al. (2023) исследовали обнаружение аномалий при distribution shift, решая проблему, когда распределения реальных данных отличаются от тех, что были видны во время обучения. Их находки предполагают, что модели, обученные на сбалансированных наборах данных, также должны быть адаптируемы к динамическим, evolving распределениям данных для поддержания высокой точности обнаружения аномалий.

Интегрируя oversampling, undersampling, синтетическую генерацию данных и distribution-aware методы обнаружения аномалий, проблемы на уровне данных в обнаружении аномалий могут быть эффективно mitigated. Эти техники усиливают robustность модели, улучшают обобщение на редкие аномалии и гарантируют лучшую применимость в реальном мире across diverse доменах, включая промышленный контроль, медицинскую диагностику и predictive maintenance.


Перспективные направления исследований
\begin{itemize}
	\item Генеративные модели с высоким качеством синтеза. Генеративные модели не должны рассматриваться лишь как инструмент аугментации данных. Если сгенерированные образцы недостаточно четкие, это окажет негативное влияние на обучение признакам. В методах контроля дефектов на основе генеративных моделей обучающая выборка часто содержит значительно больше сгенерированных образцов, чем реальных. Простое увеличение количества образцов не способно повысить итоговую точность обнаружения. Диффузионные модели (DM) в настоящее время являются наиболее передовыми и мощными генеративными моделями, обладающими большим потенциалом в задаче обнаружения дефектов. Однако эта модель все еще требует доработок из-за сложностей в обучении, низкой скорости семплирования и других причин. Важным направлением исследований будет решение проблем DM и создание методов контроля на основе этих моделей.

	\item Перенос знаний (Knowledge Transferring). Между сгенерированными и реальными образцами существуют различия в распределении. Из-за трудностей с получением новых образцов эти различия могут серьезно влиять на качество генерации и результаты контроля. Перспективным направлением исследований является извлечение общего знания и подавление специфического знания в различных наборах данных о дефектах. Например, и на стальных, и на алюминиевых листах встречаются дефекты царапин, имеющие схожие признаки. Признаки царапин являются общим знанием, в то время как различные текстуры поверхности — их специфическим знанием. Если в наборе данных по алюминиевым листам мало образцов, общее знание из наборов данных по стальным листам может быть перенесено для улучшения генерации образцов. Наиболее распространенным методом переноса знаний является кросс-доменная адаптация [126]. Эффективность извлечения знаний из синтетических данных пока неудовлетворительна, и эта исследовательская проблема требует дальнейшего изучения.

	\item Точное обнаружение или сегментация дефектов в режиме слабого контроля. Поскольку сгенерированные образцы не размечены, большинство методов контроля дефектов работают в режиме слабого контроля. Несмотря на эффективность в классификации дефектов, точность их обнаружения и сегментации значительно уступает методам с полным контролем. С одной стороны, это связано с ограничениями самого слабо контролируемого обучения, а с другой — с отсутствием надежных методов присвоения меток, что заслуживает изучения в будущих работах.

	\item Создание облегченных моделей. Существуют успешные примеры создания облегченных сверточных нейронных сетей, например, MobileNet [127]. Однако эти решения неприменимы к генеративным моделям. Для достижения высокого качества генерации мощные генеративные модели требуют сложного процесса обучения, такого как состязательное обучение в GAN или процесс диффузии и шумоподавления в DM, который зависит от громоздкой сетевой архитектуры. В итоге процесс генерации всегда требует много времени. Для обеспечения контроля в реальном времени генеративные модели требуют облегченного проектирования, что также может снизить потребность в объеме выборок.
\end{itemize} \cite{YuHe2024}
	Перспективные направления исследований
	На основе анализа раздела 7 обзора \cite{Bhatt2021} можно выделить:
	\begin{enumerate}
		\item Разработка полностью неконтролируемых методов кластеризации
		\item Интеграция с transfer learning для использования знаний из смежных доменов
		\item Создание адаптивных систем для работы с неограниченным числом типов дефектов
		\item Решение проблемы ограниченности данных через генеративные модели (GAN)
		\item Повышение объяснимости результатов кластеризации
	\end{enumerate}

\section{Цели и задачи собственного исследования}
Исходя из вышеизложенного были поставлены следующие задачи для достижения цели исследования:
\begin{enumerate}
	\item Исследовать и проанализировать существующие методы выявления аномалий изображений;
	\item Разработать метод моделирования
	\begin{enumerate}
		\item Провести сбор и первичную обработку данных;
		\item Разработать требования к входным данным;
		\item Разработать требования к выходным данным модели;
		\item Разработать метод оценки качества работы модели;
		\item Спроектировать архитектуру модели;
	\end{enumerate}
	\item Реализовать компьютерную программу модели;
	\begin{enumerate}
		\item Реализация архитектуры модели;
		\item Провести обучение модели;
	\end{enumerate}
	\item Провести тестирование модели;
	\item Провести анализ результатов работы модели.
\end{enumerate}

\section{Выводы}
В связи с этим особую значимость приобретают две задачи: повышение адаптивности традиционных методов обработки изображений для обнаружения промышленных дефектов и снижение вычислительной сложности методов глубокого обучения, применимых в условиях ограниченной разметки, а также сокращение потребности в высокопроизводительных вычислительных ресурсах \cite{Hutten2024}.

Таким образом, в результате проведения настоящего обзора были сформулированы цели и задачи будущего исследования для магистерской диссертации. В ходе работы был выполнен сбор, систематизация и анализ научной литературы, соответствующей тематике исследования, включая современные методы компьютерного зрения, алгоритмы обнаружения дефектов и подходы к обучению моделей при ограниченном количестве данных. Полученные результаты сформировали основу для разработки дальнейшей методологии исследования и определения направлений практической реализации предложенных подходов.

\printbibliography
\end{document}